\documentclass[10pt,a4paper]{amsart}
\usepackage[margin=19mm]{geometry}
\usepackage{amsmath,amssymb,mathtools}
\usepackage[scr=rsfs,cal=euler]{mathalfa}
\usepackage{xfrac}
\usepackage[unicode,hidelinks,bookmarks=false]{hyperref}
\newcommand{\ie}{i.e.\ }
\newcommand{\cf}{cf.\ }
\newcommand{\resp}{respectively\ }
\newcommand{\Subsection}{\subsection}
\providecommand{\nobreakdash}{\nobreak\hbox{-}}
\setlength{\emergencystretch}{2em}
\multlinegap0pt
\usepackage{amssymb}
\usepackage{enumerate}
\usepackage{tikz}
\newtheorem{lem}{Lemma}
\title{Stabilization index of V-number of powers of edge ideals of graphs}
\author{Ha Thi Thu Hien, Thanh Vu}
\date{Source: arXiv:2609.09632v1 (CC BY 4.0)}
\begin{document}
\maketitle

\noindent\emph{Source and licence.} Stabilization index of V-number of powers of edge ideals of graphs. Version arXiv:2609.09632v1 (CC BY 4.0), distributed by arXiv under the Creative Commons Attribution 4.0 International licence, http://creativecommons.org/licenses/by/4.0/, as recorded in the arXiv licence metadata for this exact version and shown on the abstract page https://arxiv.org/abs/2609.09632v1. Full licence notice: \texttt{licenses/arxiv-2609.09632-CC-BY-4.0.txt}.\par\smallskip

\noindent\textit{Editorial scope.} Counted unit: the article's lem lem\_admissible (source0.tex lines 505--507). The article's complete original proof of it is delivered with it (source0.tex lines 509--554, one \texttt{proof} environment of 46 lines). No mathematics is written, repaired or re-derived: the statement and the proof are the article's own lines, in the article's own notation, and no retained line is substituted or re-flowed.  Retained uncounted as the source-local context the proof needs: lines 403--408; lines 466--468.  Nothing was omitted from any retained span, so the package carries no omission record.  No retained line contains a citation, so the wrapper carries no bibliography and no bibliography entry.  No retained line contains a figure environment, an image-inclusion command or drawing code, so no image is delivered and the package declares no asset dependency.  Editorial formatting only, in addition: an AMS \texttt{amsart} preamble that restates the article's own declarations and macros used by the retained lines, namely the environments \texttt{lem} on separate counters, together with the standard packages those lines need.  The retained environments print in this document as Lemma 1 in order of appearance, whereas the article prints the same items with the numbering of its own full document.  The complete native source of the article as distributed in the arXiv e-print for this exact version is bundled unchanged as \texttt{sources/209838/source0.tex}.
\begin{lem}\label{lem_index_0} 
Let $G$ be a simple connected graph. Then 
\[
\operatorname{vstab}(I(G)) = \min \{t \ge 1 \mid v(I(G)^t) = 2t-1\}.
\]
\end{lem}

\begin{lem}\label{lem_kappa_2} 
Let $G$ be a simple connected graph. Then $\operatorname{vstab}(I(G)) \le \kappa(G)$.
\end{lem}

\begin{lem}\label{lem_admissible} 
Let $(U,V)$ be an admissible pair of $G$. Then $\operatorname{vstab}(I(G)) \le \lambda(U,V)$. 
\end{lem}

\begin{proof} 
By Lemma~\ref{lem_kappa_2}, we may assume that $U$ is non-empty. Let $U_1, \ldots, U_k$ be the connected components of $G[U]$. For each $1 \le i \le k$, let $\mathbf{a}_i \in \mathbb{N}_{>0}^{|U_i|}$ be an exponent vector such that $G[U_i]_{\mathbf{a}_i}$ is factor-critical and $\frac{|\mathbf{a}_i| - 1}{2} = \mu^*(G[U_i])$. 

Since $G^2[U \cup V]$ is connected, $V$ is an independent set, and $N_G[U] \cap V = \emptyset$, contracting each component $U_i$ to a single node $[U_i]$ yields a contracted graph $H$ where $X = \{[U_1], \ldots, [U_k]\} \cup V$ is an independent set and $H^2[X]$ is connected.

Consider a spanning tree $T$ of $H^2[X]$ rooted at $[U_1]$. For every node $v \in X \setminus \{[U_1]\}$, let $p(v)$ denote its parent in $T$. By definition of $H^2[X]$, there exists a vertex $c(v) \in V(G) \setminus (U \cup V)$ such that $c(v) \in N_G(p(v))$ and $c(v) \in N_G(v)$ (where $N_G([U_j]) = N_G(U_j)$).

Now, for each $v \in X \setminus \{[U_1]\}$, we define the monomial $m_v \in S$ as:
\[
m_v = 
\begin{cases}
x_{c(v)} x_v & \text{if } v \in V, \\
x_{c(v)} x^{\mathbf{a}_j} & \text{if } v = [U_j] \text{ for some } 1 \le j \le k.
\end{cases}
\]
Set $f = x^{\mathbf{a}_1} \prod_{v \in X \setminus \{[U_1]\}} m_v$. The degree of $f$ is given by:
\[
\deg(f) = \sum_{i=1}^k |\mathbf{a}_i| + (k - 1) + 2|V| = 2 \left( \sum_{i=1}^k \mu^*(G[U_i]) + k + |V| \right) - 1 = 2\lambda(U,V) - 1.
\]

For each $v \in X$, define $\mu^*(v) = 1$ when $v \in V$, and $\mu^*([U_j]) = \mu^*(G[U_j]) + 1$ when $v = [U_j]$. 

Notice that for every node $v \in X \setminus \{[U_1]\}$:
\begin{itemize}
    \item If $v \in V$, then $m_v = x_{c(v)} x_v \in I(G) = I(G)^{\mu^*(v)}$ since $\{c(v), v\} \in E(G)$.
    \item If $v = [U_j]$, choose $u \in U_j$ such that $c(v) \in N_G(u)$. Writing $m_v = (x_{c(v)} x_u) x^{\mathbf{a}_j'}$, the factor-criticality of $G[U_j]_{\mathbf{a}_j}$ implies $x^{\mathbf{a}_j'} \in I(G)^{\mu^*(G[U_j])}$, so $m_v \in I(G)^{1 + \mu^*(G[U_j])} = I(G)^{\mu^*([U_j])}$.
\end{itemize}

Set $t = \lambda(U,V)$. We claim that $I(G)^t : f = P_C$, where $C = N_G[U] \cup N_G(V)$ is the vertex cover of $G$ associated with the prime ideal $P_C = (x_c \mid c \in C)$. 

First, if $c \notin C$, then $c \notin N_G(\operatorname{supp}(f))$, which implies $x_c f \notin I(G)^t$. Thus, $I(G)^t : f \subseteq P_C$. 

To show the reverse inclusion $P_C \subseteq I(G)^t : f$, take any $x_c$ with $c \in C$. By definition, there exists a node $v \in X$ such that $c \in N_G(v)$ (if $v \in V$) or $c \in N_G[U_j]$ (if $v = [U_j]$). Let $[U_1] = v_0, v_1, \ldots, v_q = v$ be the path from $[U_1]$ to $v$ in $T$. We prove by induction on $q \ge 0$ that 
\[
x_c x^{\mathbf{a}_1} \prod_{i=1}^q m_{v_i} \in I(G)^s, \quad \text{where } s = \sum_{i=0}^q \mu^*(v_i).
\]
\begin{itemize}
    \item Base case ($q = 0$): Here $v = [U_1]$ and $c \in N_G[U_1]$. Since $G[U_1]$ is a connected non-bipartite component, $c$ has a neighbor $u \in U_1$ (whether $c \in U_1$ or $c \in N_G(U_1)$). Thus $x_c x_u \in I(G)$. Writing $x_c x^{\mathbf{a}_1} = (x_c x_u) x^{\mathbf{a}_1'}$, factor-criticality yields $x^{\mathbf{a}_1'} \in I(G)^{\mu^*(G[U_1])}$, whence
    \[
    x_c x^{\mathbf{a}_1} \in I(G)^{1 + \mu^*(G[U_1])} = I(G)^{\mu^*([U_1])}.
    \]

    \item Inductive step ($q \ge 1$): If $v_q \in V$, then $x_c m_{v_q} = x_c x_{c(v_q)} x_{v_q} = (x_c x_{v_q}) x_{c(v_q)} \in I(G) \cdot x_{c(v_q)}$. The remaining variable $x_{c(v_q)}$ attaches to the parent node $p(v_q) = v_{q-1}$, allowing us to apply the induction hypothesis to the path of length $q - 1$. If $v_q = [U_j]$, then $c \in N_G[U_j]$. By the same argument as the base case, $x_c x^{\mathbf{a}_j} \in I(G)^{\mu^*([U_j])}$. Factoring out these generators leaves $x_{c([U_j])}$, which attaches to $p([U_j]) = v_{q-1}$, completing the induction step.
\end{itemize}
Thus $x_c f \in I(G)^t$, so $P_C \subseteq I(G)^t : f$, giving $I(G)^t : f = P_C$. By Lemma~\ref{lem_index_0}, $\operatorname{vstab}(I(G)) \le t = \lambda(U,V)$. The conclusion follows.
\end{proof}

\end{document}
